% Develop a LaTeX script to demonstrate the presentation of Numbered theorems, definitions,
% corollaries, and lemmas in the document

\documentclass{article}

\usepackage{amsthm} % For theorem environments
\usepackage{amsmath} % For math formatting
\usepackage{lipsum} % Dummy text (optional)

% Define Theorem Styles
\newtheorem{theorem}{Theorem}[section]
\newtheorem{definition}[theorem]{Definition}
\newtheorem{corollary}[theorem]{Corollary}
\newtheorem{lemma}[theorem]{Lemma}
\newtheorem{example}[theorem]{Example}

\begin{document}

\title{Presentation of Theorems, Definitions, Corollaries, and Lemmas}
\author{Sudhanva S}
\date{\today}
\maketitle

\section{Theorems and Definitions}

\begin{theorem}
If \( a^2 + b^2 = c^2 \), then \( (a, b, c) \) is a Pythagorean triplet.
\end{theorem}

\begin{proof}
Using the Pythagorean theorem, we can verify that for values such as \( (3, 4, 5) \), the equation holds: \[ 3^2 + 4^2 = 9 + 16 = 25 = 5^2 \]
Thus, the statement is proven.
\end{proof}

\begin{definition}
A \textbf{prime number} is a natural number greater than \(1\) that has no divisors other than \(1\) and itself.
\end{definition}

\begin{example}
The numbers \(2, 3, 5, 7\) are prime numbers, while \(4\) and \(6\) are not because they have additional divisors.
\end{example}

\begin{corollary}
If \( p \) is a prime number and divides \( ab \), then \( p \) must divide either \( a \) or \( b \).
\end{corollary}

\begin{lemma}
If \( n \) is an even integer, then \( n^2 \) is also even.
\end{lemma}

\end{document}